\subsection{Frozen receivers and temporal input availability}
\label{sec:receivers}
For training mean $\boldsymbol\mu$ and per-channel standard deviation
$\boldsymbol\sigma$, let
$\mathbf z_i=(\mathbf v_i-\boldsymbol\mu)/\boldsymbol\sigma$.
Pointwise assessment uses
\begin{equation}
 s_{\mathrm{level}}(i)=\max_c|z_{i,c}|,\qquad
 s_r(i)=\frac{1}{6}\left\|\mathbf z_i-Q_rQ_r^\top\mathbf z_i\right\|_2^2,
 \quad r\in\{1,\ldots,5\},
\end{equation}
where $Q_r$ contains the leading standardized-covariance eigenvectors. Score
thresholds are the eligible-training 0.999 quantiles, and assessment uses strict
exceedance. Ranks share normalization and are correlated sensitivity analyses.
They do not constitute six independent anomaly-detector families.

The AR(1) receiver predicts a standardized observation from its original native
predecessor using a six-channel linear map and intercept. Ridge fitting minimizes
squared residuals plus $0.001n$ times the squared slope norm over $n$ eligible
training pairs; the intercept is unpenalized. Per-channel training residual RMS
scales normalize the maximum absolute forecast residual. Its threshold is also
the training 0.999 quantile. Both endpoints must be eligible for fitting, and
pairs with native gaps over 150 seconds are excluded. The first observation of a
partition has no predecessor available to the receiver.

Let $a_i$ indicate retention of original observation $i$, and let $h_i$ indicate
that its full-archive temporal score exceeds the frozen threshold. Usable temporal
evidence is
\begin{equation}
 u_i=a_{i-1}a_i h_i\,\mathbf 1\{t_i-t_{i-1}\leq150\ \mathrm{s}\}.
 \label{eq:temporal-evidence}
\end{equation}
The predecessor is the immediately preceding original observation, not merely
the preceding retained one. This explicit requirement prevents missing inputs
from being repaired using unrecorded data. It also exposes a tradeoff that a
pointwise event score alone cannot measure.
